\documentclass[11pt,a4paper]{article}
\usepackage{oi-statement}
\OISetup{CSP-J 模拟赛}{第五套}{入门级}
\OISetCompileOptions{-O2 -std=c++14}
\begin{document}
\OICoverTitle
\OICoverTable{4}{
题目名称 & 上下五千年 & 采购 & 整除序列 & 绿野仙踪\\\hline

题目类型 & 传统型 & 传统型 & 传统型 & 传统型\\\hline

目录 & {\ttfamily date} & {\ttfamily purchase} & {\ttfamily divisible} & {\ttfamily green}\\\hline

可执行文件名 & {\ttfamily date} & {\ttfamily purchase} & {\ttfamily divisible} & {\ttfamily green}\\\hline

输入文件名 & {\ttfamily date.in} & {\ttfamily purchase.in} & {\ttfamily divisible.in} & {\ttfamily green.in}\\\hline

输出文件名 & {\ttfamily date.out} & {\ttfamily purchase.out} & {\ttfamily divisible.out} & {\ttfamily green.out}\\\hline

每个测试点时限 & 1.0 秒 & 1.0 秒 & 1.0 秒 & 1.0 秒\\\hline

内存限制 & 512 MiB & 512 MiB & 512 MiB & 512 MiB\\\hline

测试点数目 & 10 & 10 & 10 & 10\\\hline

测试点是否等分 & 是，每点 10 分 & 是，每点 10 分 & 是，每点 10 分 & 是，每点 10 分\\\hline
}
\OICompilation{4}{
对于 C++ 语言 & {\ttfamily date.cpp} & {\ttfamily purchase.cpp} & {\ttfamily divisible.cpp} & {\ttfamily green.cpp}\\\hline
}
\OICoverNotes
\begin{enumerate}
\item 本套为 CSP-J 训练模拟赛，共 4 题，每题 100 分，总分 400 分。
\item 各题均采用标准输入输出，提交对应的同名 C++ 源程序。
\item 每题有 10 个测试点，每点 10 分，按测试点独立计分。
\item 答案采用普通文本比较，忽略行末空格和文件末尾多余换行。
\item 题面样例与下发大样例不计入正式测试点。
\item 程序可使用的栈空间与题目的内存限制一致。
\item 样例框左侧行号不属于输入输出。
\item 使用 C++14，编译选项为 \texttt{-O2 -std=c++14}。
\end{enumerate}

\OIProblem{上下五千年}{date}

\OISection{题目描述}

把一个日期的年份、月份和日期直接写成十进制并依次连接，中间不补前导零。若得到的字符串是回文串，则称这个日期为“回文时间”。

例如，$2018$ 年 $10$ 月 $2$ 日连接后得到 \texttt{\detokenize{2018102}}，它是回文串，所以这是一个回文时间。

日期必须是合法的公历日期：大月 $31$ 天，小月 $30$ 天；闰年二月 $29$ 天，平年二月 $28$ 天。闰年规则为：年份能被 $400$ 整除，或者能被 $4$ 整除但不能被 $100$ 整除。

给定一个合法日期，请求出\textbf{严格晚于该日期}的第一个回文时间。

\OISection{输入格式}

输入一行三个整数 $y,m,d$，分别表示年、月、日。

\OISection{输出格式}

输出一行三个整数，表示下一个回文时间的年、月、日。

\begin{OISample}{1}{输入}
2018 10 2
\end{OISample}
\begin{OISample}{1}{输出}
2019 10 2
\end{OISample}
\begin{OISample}{2}{输入}
2111 1 2
\end{OISample}
\begin{OISample}{2}{输出}
2111 1 12
\end{OISample}
\OIAdditionalSample{3}{date/date}
\OIAdditionalSample{4}{date/date}
\Needspace{9\baselineskip}
\OISection{数据范围}

\begin{itemize}
\item $2000\le y\le7000$；
\item $1\le m\le12$；
\item 输入日期合法。
\end{itemize}

本题共 $10$ 个测试点，每点 $10$ 分。

\begin{OIDataTable}{|c|c|c|}
测试点 & 年份范围 & 累计分值\\\hline
$1\sim3$ & $2000\le y\le2200$ & $30$ 分\\\hline
$1\sim6$ & $2000\le y\le4000$ & $60$ 分\\\hline
$1\sim10$ & $2000\le y\le7000$ & $100$ 分\\\hline
\end{OIDataTable}

\OIProblem{采购}{purchase}

\OISection{题目描述}

牛牛要从 $n$ 种商品中采购，每种商品只有一件。第 $i$ 件商品的体积为 $w_i$，价值为 $val_i$。所选商品总体积不能超过 $m$。

若选择商品的总价值为 $VAL$、商品数量为 $COUNT$，定义本次选择的收益为

\[
f(k)=VAL-k\times COUNT\text{。}
\]

牛牛可以一件商品都不选择。求收益的最大值。

\OISection{输入格式}

第一行包含三个正整数 $n,m,k$。

接下来 $n$ 行，每行包含两个正整数 $w_i,val_i$。

\OISection{输出格式}

输出一行一个整数，表示最大收益。

\begin{OISample}{1}{输入}
4 10 1
5 5
3 4
2 1
9 9
\end{OISample}
\begin{OISample}{1}{输出}
8
\end{OISample}
\begin{OISample}{2}{输入}
1 10 10
10 9
\end{OISample}
\begin{OISample}{2}{输出}
0
\end{OISample}
\OIAdditionalSample{3}{purchase/purchase}
\OIAdditionalSample{4}{purchase/purchase}
\Needspace{9\baselineskip}
\OISection{数据范围}

\begin{itemize}
\item $1\le n\le1000$；
\item $1\le m,k\le3000$；
\item $1\le w_i,val_i\le10^5$。
\end{itemize}

本题共 $10$ 个测试点，每点 $10$ 分，下列前两类独立计分。

\begin{OIDataTable}{|c|c|c|}
测试点 & 特殊限制 & 分值\\\hline
$1\sim2$ & $n\le20$ & $20$ 分\\\hline
$3\sim5$ & $m\le300$ & $30$ 分\\\hline
$6\sim10$ & 无额外限制 & $50$ 分\\\hline
\end{OIDataTable}

\OIProblem{整除序列}{divisible}

\OISection{题目描述}

给定两个整数 $L,R$，需要构造一个严格递增的整数序列

\[
a_1<a_2<\cdots<a_S,
\]

满足：

\begin{itemize}
\item $L\le a_i\le R$；
\item 对任意 $i<j$，都有 $a_j\bmod a_i=0$。
\end{itemize}

求能够构造出的最大长度 $S$，以及长度恰好为 $S$ 的不同序列数量。

两个序列不同，当且仅当长度不同或至少一个位置上的数不同。

本题包含多组询问，方案数需要输出完整十进制表示。

\OISection{输入格式}

第一行输入一个正整数 $T$。

接下来 $T$ 行，每行输入两个正整数 $L,R$。

\OISection{输出格式}

对每组询问输出一行两个整数：最大长度 $S$ 和达到最大长度的序列数量。

\begin{OISample}{1}{输入}
3
3 11
13 37
1 22
\end{OISample}
\begin{OISample}{1}{输出}
2 4
2 6
5 1
\end{OISample}
\OIAdditionalSample{2}{divisible/divisible}
\OIAdditionalSample{3}{divisible/divisible}
\Needspace{9\baselineskip}
\OISection{数据范围}

\begin{itemize}
\item $1\le T\le10$；
\item $1\le L<R\le10^5$；
\item 所有询问满足 $\sum R\le2\times10^5$。
\end{itemize}

本题共 $10$ 个测试点，每点 $10$ 分，下列前两类独立计分。

\begin{OIDataTable}{|c|c|c|}
测试点 & 特殊限制 & 分值\\\hline
$1\sim2$ & 对每组询问均有 $R-L\le20$ & $20$ 分\\\hline
$3\sim5$ & 对每组询问均有 $R\le2000$ & $30$ 分\\\hline
$6\sim10$ & 无额外限制 & $50$ 分\\\hline
\end{OIDataTable}

\OIProblem{绿野仙踪}{green}

\OISection{题目描述}

公园中有一块 $n\times m$ 的矩形方格场地，方格坐标满足

\[
0\le x<n,\qquad0\le y<m\text{。}
\]

公园准备以 $k$ 个给定方格 $(x_i,y_i)$ 为中心种植大小相同的正方形花圃。给定正整数 $r$，每个花圃边长为 $2r-1$，覆盖所有满足

\[
|x-x_i|\le r-1,\qquad |y-y_i|\le r-1
\]

的场内方格。超出场地的部分忽略，不同花圃的重叠方格只计算一次。

给定游园起点 $(s_x,s_y)$ 和终点 $(e_x,e_y)$。游客每一步可以向上、下、左、右移动到相邻方格，不能进入花圃覆盖的方格。花圃也不能覆盖起点或终点。

请选择统一的 $r$，使起点到终点至少存在一条合法路径，并最大化所有花圃覆盖的方格总数。若不存在任何合法的正整数 $r$，输出 $0$。

\OISection{输入格式}

第一行包含三个正整数 $n,m,k$。

接下来两行分别包含起点和终点坐标。

接下来 $k$ 行，每行包含一个花圃中心坐标 $(x_i,y_i)$。

\OISection{输出格式}

输出最大覆盖方格数。

\begin{OISample}{1}{输入}
5 6 4
0 0
2 3
0 4
2 0
2 5
4 3
\end{OISample}
\begin{OISample}{1}{输出}
21
\end{OISample}
\begin{OISample}{2}{输入}
3 5 2
0 4
2 0
0 0
2 4
\end{OISample}
\begin{OISample}{2}{输出}
8
\end{OISample}
\OIAdditionalSample{3}{green/green}
\OIAdditionalSample{4}{green/green}
\Needspace{13\baselineskip}
\OISection{数据范围}

\begin{itemize}
\item $2\le n,m\le10^3$；
\item $1\le k\le\min(nm-2,10^3)$；
\item 坐标均在场地内；
\item 花圃中心、起点、终点两两不同。
\end{itemize}

本题共 $20$ 个测试点，每点 $5$ 分。

\begin{OIDataTable}{|c|c|c|c|}
测试点 & $n,m\le$ & $k\le$ & 特殊性质\\\hline
$1\sim2$ & $50$ & $10$ & 无\\\hline
$3\sim4$ & $1000$ & $1$ & 无\\\hline
$5\sim6$ & $1000$ & $10$ & 无\\\hline
$7\sim8$ & $1000$ & $1000$ & 起点和终点相邻\\\hline
$9\sim10$ & $1000$ & $1000$ & 起点为 $(0,0)$，终点为 $(n-1,m-1)$\\\hline
$11\sim20$ & $1000$ & $1000$ & 无\\\hline
\end{OIDataTable}

\end{document}
